% from main.tex
%%%%%%%%%%  scan idx 145  |  folio 132  %%%%%%%%%%
\textbf{Définition 5.2.} Un objet \(F\) de \(\Dpx^-({}_A S)\) satisfaisant aux conditions équivalentes de (5.1) sera dit d'amplitude plate contenue dans l'intervalle \([m,n]\), et nous écrirons
\[
  \toramp(F)\subset[m,n].
\]
Nous dirons que \(F\) est d'amplitude plate finie (ou de tor-dimension finie) s'il existe des entiers \(m\) et \(n\) tels que la condition précédente soit vérifiée.

\textbf{5.3.} Soit \(X\in\Ob S\). Nous noterons \(\Dpx^-({}_A X)_{\torf}\) la sous-catégorie pleine de \(\Dpx^-({}_A X)\) formée des objets de tor-dimension finie. D'après le critère (ii) de (5.1), c'est une sous-catégorie triangulée de \(\Dpx^-({}_A X)\). Pour \(X\) variable, les catégories \(\Dpx^-({}_A X)_{\torf}\) forment une sous-\(S\)-catégorie triangulée (1.1.4) de \(\Dpx^-({}_A S)\), que nous noterons \(\Dpx^-({}_A S)_{\torf}\). Le foncteur \(\overset{\mathrm L}{\otimes}_A\) induit un \(S\)-foncteur exact
\[
  \overset{\mathrm L}{\otimes}_A:
  \Dpx^b(S_A)\times_S \Dpx({}_A S)_{\torf}
  \longrightarrow \Dpx^b(S_{\mathbb Z}).
\]

\textbf{5.4.} Pour qu'un \(A\)-Module à gauche \(F\) soit d'amplitude plate contenue dans \([-n,0]\) (\(n\in\mathbb N\)) il faut et il suffit que \(F\) soit de tor-dimension \(\leq n\) au sens ordinaire, i.e. que \(F\) admette une résolution gauche de longueur \(n\) par des \(A\)-Modules plats.

\textbf{5.5.}
